\section{Theoretical Analysis}
\label{sec:theory-main}

The following properties specify the architecture's retention and interface
contracts. They concern exact finite semantics and real-arithmetic sufficient
conditions; numerical execution additionally requires its declared precision
and clearance checks. Detailed constructions and proofs appear in the appendix.

\paragraph{Operation preservation.}
Let $\delta$ be an admitted polygonal loop's minimum distance from the
discriminant. If corresponding vertices move by less than $\delta$, every point
of their linear interpolation moves by at most that amount. The interpolation
therefore avoids the discriminant and preserves the endpoint permutation.
Preserving each generator preserves every finite composition under fixed state
ports. This property permits implementation changes while retaining the
acquired operation.

\paragraph{Task and soft-interface preservation.}
For changing neural bindings, require a finite relation between reference and
candidate states that contains every old encoded input pair, agrees on decoded
outputs, and is closed under every old token transition. Induction on program
length then preserves all native old words. Separately, let $p_E,p_R,p_C$ be
minimum probabilities assigned to the model's own hard encoding, readout, and
call choices. With $S$ slots and one-hot executed maps,
\begin{equation}
 m_L=p_Ep_Rp_C^{SL}>\tfrac12
 \label{eq:main-path-mass}
\end{equation}
implies soft/native argmax agreement through length $L$: the designated path
alone carries more probability than all alternatives combined. The native
closure and finite-length soft guarantee have different scopes.

\paragraph{Irreversible logical interfaces.}
Let $Y$ be an $N$-element logical state set, let $f:Y\to Y$, and suppose
its marked frame is a section of a physical decoder. Let
$G$ be its retained permutation group, and require old operations to remain
decoder-equivariant on an invariant set containing both the frame and the
outputs of one reversible lift of $f$. The minimum physical-state count is
\begin{equation}
 D_{1,G}=\sum_{O\in Y/G}|O|
       \max\!\left(1,\max_{y\in O}|f^{-1}(y)|\right).
\end{equation}
Fibers have equal cardinality along each orbit and must accommodate distinct
preimages; assigning sufficient copies and extending a partial injection proves
attainability. For transitive $G$, reset requires $N^2$ states. Repeated
non-surjective reset cannot be a closed finite reversible quotient: injectivity
makes its forward-invariant reachable set invariant, forcing the logical map
to be surjective. Fresh ancillas or an explicit irreversible interface change
this contract and incur storage or disposal costs. Configuration counts remain
distinct from factorized memory bits and trainable parameters.
